\begin{ex}
Tính: $\frac{1}{4} + \frac{2}{4} = \dots$
\shortans{3/4}
\loigiai{
$\frac{3}{4}$.
}
\end{ex}

\begin{ex}
Tính: $\frac{5}{6} - \frac{1}{6} = \dots$ (tối giản).
\shortans{2/3}
\loigiai{
$\frac{4}{6} = \frac{2}{3}$.
}
\end{ex}

\begin{ex}
Tính: $\frac{1}{2} + \frac{1}{3} = \dots$
\shortans{5/6}
\loigiai{
$\frac{5}{6}$.
}
\end{ex}

\begin{ex}
Tính: $\frac{1}{2} - \frac{1}{3} = \dots$
\shortans{1/6}
\loigiai{
$\frac{1}{6}$.
}
\end{ex}

\begin{ex}
Tính: $\frac{2}{5} + \frac{1}{10} = \dots$ (tối giản).
\shortans{1/2}
\loigiai{
$\frac{4}{10} + \frac{1}{10} = \frac{5}{10} = \frac{1}{2}$.
}
\end{ex}

\begin{ex}
Tính: $\frac{7}{8} - \frac{1}{4} = \dots$
\shortans{5/8}
\loigiai{
$\frac{7}{8} - \frac{2}{8} = \frac{5}{8}$.
}
\end{ex}

\begin{ex}
Tính: $1 - \frac{4}{7} = \dots$
\shortans{3/7}
\loigiai{
$\frac{3}{7}$.
}
\end{ex}

\begin{ex}
Tính: $1 + \frac{3}{5} = \dots$
\shortans{8/5}
\loigiai{
$\frac{8}{5}$.
}
\end{ex}

\begin{ex}
Tính giá trị biểu thức: $\frac{1}{3} + \frac{1}{4} + \frac{5}{12}$ (kết quả bằng số tự nhiên).
\shortans{1}
\loigiai{
$\frac{4}{12} + \frac{3}{12} + \frac{5}{12} = \frac{12}{12} = 1$.
}
\end{ex}

\begin{ex}
Tính giá trị biểu thức: $1 - \frac{1}{3} - \frac{1}{3}$.
\shortans{1/3}
\loigiai{
$1 - \frac{2}{3} = \frac{1}{3}$.
}
\end{ex}

\begin{ex}
Tìm $x$: $x + \frac{1}{5} = \frac{4}{5}$.
\shortans{3/5}
\loigiai{
$x = \frac{4}{5} - \frac{1}{5} = \frac{3}{5}$.
}
\end{ex}

\begin{ex}
Tìm $x$: $x - \frac{1}{2} = \frac{1}{4}$.
\shortans{3/4}
\loigiai{
$x = \frac{1}{4} + \frac{2}{4} = \frac{3}{4}$.
}
\end{ex}

\begin{ex}
Tìm $x$: $\frac{7}{9} - x = \frac{2}{9}$.
\shortans{5/9}
\loigiai{
$x = \frac{7}{9} - \frac{2}{9} = \frac{5}{9}$.
}
\end{ex}

\begin{ex}
Một mảnh đất được chia làm 3 phần: phần trồng hoa chiếm $\frac{1}{4}$, phần làm lối đi chiếm $\frac{1}{4}$. Phần còn lại để trồng rau chiếm bao nhiêu phần mảnh đất?
\shortans{1/2}
\loigiai{
$1 - (\frac{1}{4} + \frac{1}{4}) = 1 - \frac{1}{2} = \frac{1}{2}$.
}
\end{ex}

\begin{ex}
Tính: $\frac{3}{8} + \frac{5}{8} = \dots$ (kết quả số tự nhiên).
\shortans{1}
\loigiai{
$1$.
}
\end{ex}

\begin{ex}
Tính: $\frac{7}{6} - \frac{1}{6} = \dots$ (kết quả số tự nhiên).
\shortans{1}
\loigiai{
$\frac{6}{6} = 1$.
}
\end{ex}

\begin{ex}
Tính nhanh: $\frac{2}{9} + \frac{5}{11} + \frac{7}{9}$.
\shortans{16/11}
\loigiai{
$(\frac{2}{9} + \frac{7}{9}) + \frac{5}{11} = 1 + \frac{5}{11} = \frac{16}{11}$.
}
\end{ex}

\begin{ex}
Tính nhanh: $\frac{3}{5} + \frac{1}{4} + \frac{2}{5}$.
\shortans{5/4}
\loigiai{
$(\frac{3}{5} + \frac{2}{5}) + \frac{1}{4} = 1 + \frac{1}{4} = \frac{5}{4}$.
}
\end{ex}

\begin{ex}
Một bể đã có $\frac{1}{3}$ lượng nước, vòi chảy thêm $\frac{1}{2}$ lượng nước nữa. Bể có bao nhiêu phần nước?
\shortans{5/6}
\loigiai{
$\frac{1}{3} + \frac{1}{2} = \frac{5}{6}$ bể.
}
\end{ex}

\begin{ex}
Lượng nước còn thiếu để đầy bể trong bài toán trên là bao nhiêu phần?
\shortans{1/6}
\loigiai{
$1 - \frac{5}{6} = \frac{1}{6}$ bể.
}
\end{ex}

\begin{ex}
Tính: $2 - \frac{1}{3} = \dots$
\shortans{5/3}
\loigiai{
$\frac{6}{3} - \frac{1}{3} = \frac{5}{3}$.
}
\end{ex}

\begin{ex}
Tính: $3 + \frac{1}{2} = \dots$
\shortans{7/2}
\loigiai{
$\frac{6}{2} + \frac{1}{2} = \frac{7}{2}$.
}
\end{ex}

\begin{ex}
Tính: $\frac{4}{7} + \frac{3}{7} = \dots$ (kết quả số tự nhiên).
\shortans{1}
\loigiai{
$\frac{7}{7} = 1$.
}
\end{ex}

\begin{ex}
Tính: $\frac{5}{12} - \frac{1}{6} = \dots$
\shortans{1/4}
\loigiai{
$\frac{5}{12} - \frac{2}{12} = \frac{3}{12} = \frac{1}{4}$.
}
\end{ex}

\begin{ex}
Tính: $\frac{2}{3} + \frac{1}{6} - \frac{1}{2} = \dots$
\shortans{1/3}
\loigiai{
$\frac{4}{6} + \frac{1}{6} - \frac{3}{6} = \frac{2}{6} = \frac{1}{3}$.
}
\end{ex}
